% ============================================================
% Enhancing Multimodal Medical Image Fusion
% Master's Thesis — Benchmark & Metrics Reference
% ============================================================
\documentclass[a4paper,12pt]{article}

\usepackage[margin=2.5cm]{geometry}
\usepackage{graphicx}
\usepackage{booktabs}
\usepackage{multirow}
\usepackage{makecell}
\usepackage{caption}
\usepackage{array}
\usepackage{amsmath, amssymb}
\usepackage{enumitem}
\usepackage{float}
\usepackage[hidelinks]{hyperref}

\newcommand{\best}[1]{\textbf{#1}}
\newcommand{\rot}[1]{\makecell{#1}}

\begin{document}

% ── Title ──
\begin{center}
{\LARGE\bfseries Enhancing Multimodal Medical Image Fusion:\\[0.3em]
Benchmark Analysis and Model Improvement}\\[1.2em]
{\large Master's Thesis}
\end{center}

\tableofcontents
\newpage

% ════════════════════════════════════════════════════════════════
\section{Evaluation Metrics}
% ════════════════════════════════════════════════════════════════

Image fusion combines two source images $A$ (e.g., CT, PET, or SPECT) and $B$ (MRI)
into a single fused image $F$ that retains complementary information from both modalities.
Since no ground-truth fused image exists, all evaluation metrics are \emph{no-reference}:
they measure how well $F$ preserves information from $A$ and $B$.

We adopt 16 metrics organized into 6 groups. The first 4 metrics follow
Xie et al.~\cite{fusionmamba2024}; the remaining 12 follow the taxonomy of
Liu et al.~\cite{liu2012}.
All metrics are reported as \textbf{higher is better} ($\uparrow$).

\medskip
\noindent\textbf{Notation.}
$A(i,j)$, $B(i,j)$, $F(i,j)$: pixel intensity at position $(i,j)$.
$M \times N$: image size.
$H(X) = -\sum_k p_k \log_2 p_k$: Shannon entropy.
$I(X;Y) = H(X) + H(Y) - H(X,Y)$: mutual information.
$\mu_X$, $\sigma_X$, $\sigma_{XY}$: local mean, standard deviation, and covariance
computed over a sliding window of size $W \times W$.

% ──────────────────────────────────────
\subsection{Structural Similarity Metrics}
% ──────────────────────────────────────

\subsubsection{MS-SSIM --- Multi-Scale Structural Similarity}

Wang et al.~\cite{wang2003} extended SSIM to multiple scales by iteratively
low-pass filtering and downsampling the images. At each scale $j$, the contrast
comparison $c_j$ and structure comparison $s_j$ are computed; the luminance
comparison $l_M$ is evaluated only at the coarsest scale $M$:
\begin{equation}
  \text{MS-SSIM}(X,Y) = l_M^{\alpha_M} \cdot \prod_{j=1}^{M} c_j^{\beta_j} \cdot s_j^{\gamma_j},
\end{equation}
where $\alpha_M = \beta_j = \gamma_j$ are set by psychophysical experiments
($M=5$, weights: 0.0448, 0.2856, 0.3001, 0.2363, 0.1333).
For fusion evaluation, we report
$\text{MS-SSIM} = \frac{1}{2}[\text{MS-SSIM}(A,F) + \text{MS-SSIM}(B,F)]$.

\textbf{Range:} $[0,1]$. A value of 1 indicates perfect structural preservation
across all scales.

\subsubsection{FMI --- Feature Mutual Information}

Haghighat et al.~\cite{haghighat2011} proposed measuring the mutual information
between feature representations (rather than raw pixels) of the source and fused images.
Features are extracted via gradient operators (Sobel), and mutual information
is computed within non-overlapping $w \times w$ blocks:
\begin{equation}
  \text{FMI} = \frac{1}{N_b} \sum_{b=1}^{N_b}
    \bigl[ I(f_A^b; f_F^b) + I(f_B^b; f_F^b) \bigr],
\end{equation}
where $f_X^b$ is the feature vector of image $X$ in block $b$, and $N_b$ is the
number of blocks.

\textbf{Range:} $[0, +\infty)$, typical: $0.18$--$0.47$.
Higher values indicate better feature-level information preservation.

% ──────────────────────────────────────
\subsection{Visual Fidelity Metrics}
% ──────────────────────────────────────

\subsubsection{VIF --- Visual Information Fidelity}

Sheikh and Bovik~\cite{sheikh2006} model the image source as a Gaussian scale
mixture (GSM) in the wavelet domain and define VIF as the ratio of information
the human visual system (HVS) extracts from the distorted image versus the
reference:
\begin{equation}
  \text{VIF}(X,F) = \frac{\sum_j \sum_k \log_2\!\bigl(1 + g_j^2 s_{jk}^2 C_u / (\sigma_v^2 + \sigma_n^2)\bigr)}
                         {\sum_j \sum_k \log_2\!\bigl(1 + s_{jk}^2 C_u / \sigma_n^2\bigr)},
\end{equation}
where $j$ indexes wavelet sub-bands, $k$ indexes spatial locations, $g_j$ and
$\sigma_v^2$ are the gain and noise of the distortion channel, $s_{jk}$ is the
local signal variance, and $\sigma_n^2$ models HVS internal noise.
For fusion: $\text{VIF} = \frac{1}{2}[\text{VIF}(A,F) + \text{VIF}(B,F)]$.

\textbf{Range:} $[0, 1+]$. VIF = 1 means all visual information is preserved;
VIF $> 1$ implies enhancement.

\subsubsection{SCD --- Sum of Correlations of Differences}

Aslantas and Bendes~\cite{aslantas2015} define SCD based on the intuition that
if $F$ successfully fuses $A$ and $B$, then the residual $F - B$ should resemble
$A$, and vice versa:
\begin{equation}
  \text{SCD} = r(F - B,\; A) + r(F - A,\; B),
\end{equation}
where $r(\cdot,\cdot)$ denotes the Pearson correlation coefficient.

\textbf{Range:} $[-2, 2]$.
SCD $= 2$ indicates perfect content separation; higher values indicate
that $F$ faithfully preserves the unique content of each source.

% ──────────────────────────────────────
\subsection{Information Theory-Based Metrics}
% ──────────────────────────────────────

The following three metrics, from Group~1 of Liu et al.~\cite{liu2012},
quantify how much statistical information is transferred from the sources to
the fused image.

\subsubsection{$Q_{MI}$ --- Normalized Mutual Information}

Hossny et al.~\cite{hossny2008} normalize MI to avoid scale dependence:
\begin{equation}
  Q_{MI} = 2\!\left[\frac{I(A,F)}{H(A)+H(F)} + \frac{I(B,F)}{H(B)+H(F)}\right].
\end{equation}

\textbf{Range:} $[0,2]$, typical: $0.5$--$0.9$ for medical images.
$Q_{MI} = 2$ means $F$ captures all information from both sources.

\subsubsection{$Q_{TE}$ --- Tsallis Entropy-Based Metric}

Cvejic et al.~\cite{cvejic2006} and Nava et al.~\cite{nava2007} use Tsallis
entropy ($q = 1.85$) as a generalization of Shannon entropy:
\begin{equation}
  Q_{TE} = \frac{I^q(A,F) + I^q(B,F)}{H^q(A) + H^q(B) - I^q(A,B)},
\end{equation}
where $H^q(X) = \frac{1}{q-1}\!\left(1 - \sum_k p_k^q\right)$ is the Tsallis
entropy and $I^q(X,Y)$ is the corresponding mutual information.

\textbf{Range:} $(-\infty, +\infty)$, unbounded.
The denominator can become negative when sources are highly correlated,
producing negative values. Best used for \emph{relative} comparison within the
same dataset.

\subsubsection{$Q_{NCIE}$ --- Nonlinear Correlation Information Entropy}

The nonlinear correlation coefficient (NCC) between images $X$ and $Y$ is defined
as $\text{NCC}(X,Y) = H'(X) + H'(Y) - H'(X,Y)$ where $H'$ uses base-$b$
logarithms ($b=256$). A $3 \times 3$ correlation matrix $R$ is formed from the
pairwise NCCs of $\{A, B, F\}$, and its eigenvalues $\lambda_1, \lambda_2, \lambda_3$
yield:
\begin{equation}
  Q_{NCIE} = 1 + \sum_{i=1}^{3} \frac{\lambda_i}{3} \log_3 \frac{\lambda_i}{3}.
\end{equation}

\textbf{Range:} $[0,1]$, typical: $0.02$--$0.04$.
Higher values indicate that $F$ participates uniformly in the correlation
structure with both sources.

% ──────────────────────────────────────
\subsection{Image Feature-Based Metrics}
% ──────────────────────────────────────

Group~2 of Liu et al.~\cite{liu2012} evaluates whether edges, textures, and
spatial details of $A$ and $B$ are preserved in $F$.

\subsubsection{$Q^{AB/F}$ --- Gradient-Based Quality}

Xydeas and Petrovic~\cite{xydeas2000} compute Sobel gradient strength
$g_X(i,j)$ and orientation $\alpha_X(i,j)$ for each image. The relative
strength $G^{AF}$ and orientation preservation $\Delta^{AF}$ are passed through
sigmoid functions to obtain edge information preservation $Q^{AF}(i,j)$.
The final metric is:
\begin{equation}
  Q^{AB/F} = \frac{\sum_{i,j}\bigl[Q^{AF}(i,j) \cdot g_A(i,j) + Q^{BF}(i,j) \cdot g_B(i,j)\bigr]}
             {\sum_{i,j}\bigl[g_A(i,j) + g_B(i,j)\bigr]}.
\end{equation}

\textbf{Range:} $[0,1]$, typical: $0.3$--$0.75$.
This is the most widely adopted edge-preservation metric. Pixels with stronger
edges receive higher weight.

\subsubsection{$Q_M$ --- Multi-Scale Wavelet Metric}

Wang and Liu~\cite{wang2008} use a two-level Haar wavelet decomposition.
At each scale $s$, the edge preservation $EP_s^{AF}$ is computed from the
high-frequency sub-bands (LH, HL, HH), weighted by the source energy:
\begin{equation}
  Q_M = \prod_{s=1}^{S} \left(Q_s^{AB/F}\right)^{\alpha_s}, \quad
  Q_s^{AB/F} = \frac{\sum_m \bigl[EP_s^{AF}(m) w_A^s(m) + EP_s^{BF}(m) w_B^s(m)\bigr]}
                     {\sum_m \bigl[w_A^s(m) + w_B^s(m)\bigr]}.
\end{equation}

\textbf{Range:} $[0,1]$.
Captures both fine-scale and coarse-scale feature preservation.

\subsubsection{$Q_{SF}$ --- Spatial Frequency Ratio}

Zheng et al.~\cite{zheng2007} define spatial frequency as
$SF = \sqrt{RF^2 + CF^2 + MDF^2 + SDF^2}$, where RF, CF, MDF, SDF are
first-order gradient energies in four directions. The metric compares
$SF(F)$ against a reference computed from the maximum gradients of $A$
and $B$:
\begin{equation}
  Q_{SF} = \frac{SF(F) - SF_{\text{ref}}}{SF_{\text{ref}}}.
\end{equation}

\textbf{Range:} $(-1, +\infty)$, typical: $-0.45$ to $0$.
$Q_{SF} = 0$: $F$ matches the best source detail level;
$Q_{SF} < 0$: over-smoothing; $Q_{SF} > 0$: enhancement.

\subsubsection{$Q_P$ --- Phase Congruency-Based Metric}

Zhao et al.~\cite{zhao2007} extract three feature maps (gradient magnitude,
Laplacian, orientation) and compute the maximum cross-correlation between
source features and fused features:
\begin{equation}
  Q_P = \rho(g_A, g_B, g_F) \cdot \rho(\ell_A, \ell_B, \ell_F) \cdot \rho(\phi_A, \phi_B, \phi_F),
\end{equation}
where $\rho$ takes the maximum correlation among $(A,F)$, $(B,F)$, and $(A+B,F)$.

\textbf{Range:} $[0,1]$, typical: $0.05$--$0.35$.
Higher values indicate better preservation of perceptually salient features
(edges, corners, ridges).

% ──────────────────────────────────────
\subsection{Structural Similarity-Based Metrics}
% ──────────────────────────────────────

Group~3 of Liu et al.~\cite{liu2012} uses local statistics (mean, variance,
covariance) to assess structural similarity. All three metrics build on the
Universal Image Quality Index (UIQI):
\begin{equation}
  Q_0(X,Y) = \frac{2\mu_X \mu_Y}{\mu_X^2 + \mu_Y^2} \cdot
              \frac{2\sigma_X \sigma_Y}{\sigma_X^2 + \sigma_Y^2} \cdot
              \frac{\sigma_{XY}}{\sigma_X \sigma_Y}.
\end{equation}

\subsubsection{$Q_S$ --- Piella's Saliency-Weighted Metric}

Piella and Heijmans~\cite{piella2003} weight the contributions of $A$ and $B$
by their local saliency (variance):
\begin{equation}
  Q_S = \frac{1}{|W|} \sum_{w \in W}
    \bigl[\lambda(w) \, Q_0(A, F|w) + (1 - \lambda(w)) \, Q_0(B, F|w)\bigr], \quad
  \lambda(w) = \frac{\sigma_A^2(w)}{\sigma_A^2(w) + \sigma_B^2(w)}.
\end{equation}

\textbf{Range:} $[-1, 1]$, practically $[0, 1]$.
The more active source at each location is weighted more.

\subsubsection{$Q_C$ --- Cvejic's Covariance-Weighted Metric}

Cvejic et al.~\cite{cvejic2005} replace the saliency weight with the
covariance between each source and $F$:
\begin{equation}
  Q_C = \frac{1}{|W|} \sum_{w} \bigl[\text{sim}(w) \, Q_0(A,F|w) + (1 - \text{sim}(w)) \, Q_0(B,F|w)\bigr], \quad
  \text{sim}(w) = \frac{\sigma_{AF}(w)}{\sigma_{AF}(w) + \sigma_{BF}(w)}.
\end{equation}

\textbf{Range:} $[-1, 1]$, practically $[0, 1]$.
Rewards fusion outputs that adaptively draw from the most correlated source.

\subsubsection{$Q_Y$ --- Yang's SSIM-Based Adaptive Metric}

Yang et al.~\cite{yang2008} use SSIM with an adaptive threshold: if sources
are similar at a location, $F$ must match both; otherwise, it suffices to
match the better source:
\begin{equation}
  Q_Y(w) = \begin{cases}
    \lambda(w) \cdot \text{SSIM}(A,F|w) + (1-\lambda(w)) \cdot \text{SSIM}(B,F|w),
      & \text{SSIM}(A,B|w) \geq 0.75, \\
    \max\!\bigl(\text{SSIM}(A,F|w),\; \text{SSIM}(B,F|w)\bigr),
      & \text{SSIM}(A,B|w) < 0.75.
  \end{cases}
\end{equation}

\textbf{Range:} $[-1, 1]$, practically $[0, 1]$.
The threshold $0.75$ makes this metric particularly suited for multimodal
medical fusion where sources are very different.

% ──────────────────────────────────────
\subsection{Human Perception-Inspired Metrics}
% ──────────────────────────────────────

Group~4 of Liu et al.~\cite{liu2012} incorporates models of the human visual
system (HVS), including contrast sensitivity filtering (CSF) and edge saliency.

\subsubsection{$Q_{CV}$ --- Chen-Varshney Perceptual Distortion Metric}

Chen and Varshney~\cite{chen2007} compute a saliency-weighted perceptual
distortion. Edge saliency $\Lambda_X(b) = \sum_{(i,j) \in b} g_X(i,j)^\alpha$
is computed per block $b$ ($\alpha = 5$), and the distortion $D_X(b)$ is
measured after contrast sensitivity filtering (Mannos--Sakrison CSF):
\begin{equation}
  Q_{CV} = -\frac{\sum_{b}\bigl[\Lambda_A(b) D_A(b) + \Lambda_B(b) D_B(b)\bigr]}
                  {\sum_{b}\bigl[\Lambda_A(b) + \Lambda_B(b)\bigr]}.
\end{equation}
The sign is negated so that higher (closer to 0) is better.

\textbf{Range:} $(-\infty, 0]$, typical: $-3900$ to $-55$.
Values near zero indicate low perceptual distortion. Large negative values
arise when $F$ has significantly different intensity distributions from the
sources in salient (edge-rich) regions.

\subsubsection{$Q_{CB}$ --- Chen-Blum Perceptual Contrast Metric}

Chen and Blum~\cite{chen2009} model the eye's band-pass contrast perception
using a Difference-of-Gaussians (DoG) CSF filter. Local contrast is computed
as the ratio of two Gaussian-filtered images, and a contrast masking model
determines the information preservation at each location:
\begin{equation}
  Q_{CB} = \frac{\sum_{i,j}\bigl[w_A(i,j) \, Q_{CM}(A,F) + w_B(i,j) \, Q_{CM}(B,F)\bigr]}
                 {\sum_{i,j}\bigl[w_A(i,j) + w_B(i,j)\bigr]},
\end{equation}
where $Q_{CM}(X,F)$ compares the masked local contrast maps, and $w_X$
is a saliency weight based on contrast magnitude.

\textbf{Range:} $[0,1]$, typical: $0.3$--$0.77$.
Higher values indicate better perceptual contrast preservation.

% ──────────────────────────────────────
\subsection{Summary}
% ──────────────────────────────────────

Table~\ref{tab:metrics-summary} summarizes all 16 metrics with their ranges
and sources. No single metric suffices: each captures different failure modes
(blur, modality bias, structural distortion, contrast loss). A comprehensive
evaluation requires all 16 metrics.

\begin{table}[H]
\centering
\caption{Summary of 16 evaluation metrics.}
\label{tab:metrics-summary}
\renewcommand{\arraystretch}{1.25}
\small
\begin{tabular}{llll}
\toprule
\textbf{Metric} & \textbf{Group} & \textbf{Range} & \textbf{Source} \\
\midrule
MS-SSIM    & Similarity & $[0,1]$              & Wang et al.~\cite{wang2003} \\
FMI        & Similarity & $[0,+\infty)$        & Haghighat et al.~\cite{haghighat2011} \\
\midrule
VIF        & Fidelity   & $[0,1+]$             & Sheikh \& Bovik~\cite{sheikh2006} \\
SCD        & Fidelity   & $[-2,2]$             & Aslantas \& Bendes~\cite{aslantas2015} \\
\midrule
$Q_{MI}$   & Info Theory    & $[0,2]$              & Hossny et al.~\cite{hossny2008} \\
$Q_{TE}$   & Info Theory    & $(-\infty,+\infty)$  & Cvejic et al.~\cite{cvejic2006} \\
$Q_{NCIE}$ & Info Theory    & $[0,1]$              & Liu et al.~\cite{liu2012} \\
\midrule
$Q^{AB/F}$ & Image Feature  & $[0,1]$              & Xydeas \& Petrovic~\cite{xydeas2000} \\
$Q_M$      & Image Feature  & $[0,1]$              & Wang \& Liu~\cite{wang2008} \\
$Q_{SF}$   & Image Feature  & $(-1,+\infty)$       & Zheng et al.~\cite{zheng2007} \\
$Q_P$      & Image Feature  & $[0,1]$              & Zhao et al.~\cite{zhao2007} \\
\midrule
$Q_S$      & Struct Sim     & $[-1,1]$    & Piella \& Heijmans~\cite{piella2003} \\
$Q_C$      & Struct Sim     & $[-1,1]$    & Cvejic et al.~\cite{cvejic2005} \\
$Q_Y$      & Struct Sim     & $[-1,1]$    & Yang et al.~\cite{yang2008} \\
\midrule
$Q_{CV}$   & Perception     & $(-\infty,0]$  & Chen \& Varshney~\cite{chen2007} \\
$Q_{CB}$   & Perception     & $[0,1]$        & Chen \& Blum~\cite{chen2009} \\
\bottomrule
\end{tabular}
\end{table}

\newpage

% ════════════════════════════════════════════════════════════════
\section{Experimental Results}
% ════════════════════════════════════════════════════════════════

We evaluate 11 fusion methods on three medical image fusion tasks from the
Harvard public medical dataset: CT-MRI (24 pairs), PET-MRI (24 pairs), and
SPECT-MRI (24 pairs). All images are $256 \times 256$ pixels. For color
source images (PET, SPECT), metrics are computed on the Y channel of the
YCbCr color space. Results are reported as the average over all test pairs.

\medskip
\noindent\textbf{Compared methods.}
CDDFuse~\cite{cddfuse} (IVF and MIF variants),
MambaDFuse~\cite{mambadfuse},
SwinFusion~\cite{swinfusion},
LRFNet~\cite{lrfnet},
IFCNN~\cite{ifcnn},
EMMA~\cite{emma},
DDFM~\cite{ddfm},
MUFusion~\cite{mufusion},
NestFuse~\cite{nestfuse},
DIDFuse~\cite{didfuse}.

\clearpage
\newgeometry{margin=1cm}

% ────────────────────────────────────────────────────────────────
%  TABLE — CT-MRI
% ────────────────────────────────────────────────────────────────
\begin{table}[H]
\centering
\caption{Quantitative comparison on \textbf{CT-MRI} (24 pairs).}
\label{tab:ctmri}
\setlength{\tabcolsep}{2.5pt}
\renewcommand{\arraystretch}{1.1}
\scriptsize
\resizebox{\textwidth}{!}{%
\begin{tabular}{l | cc | cc | ccc | cccc | ccc | cc | cc}
\toprule
 &
\multicolumn{2}{c|}{\textbf{Similarity}} &
\multicolumn{2}{c|}{\textbf{Fidelity}} &
\multicolumn{3}{c|}{\textbf{Info Theory}} &
\multicolumn{4}{c|}{\textbf{Image Feature}} &
\multicolumn{3}{c|}{\textbf{Struct Sim}} &
\multicolumn{2}{c|}{\textbf{Perception}} &
\\
\cmidrule(lr){2-3}\cmidrule(lr){4-5}\cmidrule(lr){6-8}\cmidrule(lr){9-12}\cmidrule(lr){13-15}\cmidrule(lr){16-17}
\textbf{Method} &
\rot{MS-SSIM} & \rot{FMI} & \rot{VIF} & \rot{SCD} &
\rot{$Q_{MI}$} & \rot{$Q_{TE}$} & \rot{$Q_{NCIE}$} &
\rot{$Q^{AB/F}$} & \rot{$Q_M$} & \rot{$Q_{SF}$} & \rot{$Q_P$} &
\rot{$Q_S$} & \rot{$Q_C$} & \rot{$Q_Y$} &
\rot{$Q_{CV}$} & \rot{$Q_{CB}$} & \rot{EN} & \rot{PSNR} \\
\midrule
CDDFuse\textsubscript{IVF} & \best{0.652} & 0.297 & 0.250 & \best{1.742} & 0.731 & $-$7.47 & 0.029 & 0.589 & 0.631 & $-$0.133 & 0.114 & 0.395 & 0.419 & 0.907 & $-$1137 & 0.635 & 4.73 & 13.55 \\
CDDFuse\textsubscript{MIF} & 0.643 & \best{0.356} & \best{0.303} & 1.400 & 0.842 & $-$12.01 & \best{0.040} & \best{0.675} & \best{0.652} & \best{$-$0.025} & \best{0.216} & \best{0.442} & \best{0.491} & 0.943 & \best{$-$1028} & \best{0.702} & 4.77 & 15.03 \\
MambaDFuse & 0.591 & 0.312 & 0.264 & 1.599 & 0.769 & 50.44 & 0.025 & 0.592 & 0.639 & $-$0.150 & 0.140 & 0.398 & 0.449 & 0.927 & $-$1236 & 0.658 & 4.60 & 13.91 \\
SwinFusion & 0.584 & 0.301 & 0.271 & 1.529 & 0.790 & 56.72 & 0.026 & 0.555 & 0.643 & $-$0.209 & 0.128 & 0.383 & 0.437 & 0.934 & $-$1212 & 0.656 & 4.54 & 14.13 \\
LRFNet & 0.607 & 0.315 & 0.249 & 1.347 & 0.778 & 53.08 & 0.026 & 0.483 & 0.640 & $-$0.403 & 0.190 & 0.335 & 0.397 & 0.894 & $-$3927 & 0.662 & 4.49 & \best{16.55} \\
IFCNN & 0.597 & 0.294 & 0.221 & 1.365 & 0.716 & 41.37 & 0.022 & 0.571 & 0.622 & $-$0.125 & 0.118 & 0.399 & 0.431 & 0.902 & $-$1215 & 0.663 & 4.60 & 15.30 \\
EMMA & 0.591 & 0.251 & 0.233 & 1.518 & 0.682 & 50.31 & 0.024 & 0.510 & 0.581 & $-$0.332 & 0.096 & 0.372 & 0.416 & 0.885 & $-$1598 & 0.551 & \best{5.36} & 14.29 \\
DDFM & 0.609 & 0.317 & 0.204 & 1.109 & \best{0.852} & \best{75.44} & 0.030 & 0.373 & 0.635 & $-$0.481 & 0.148 & 0.338 & 0.353 & 0.834 & $-$2042 & 0.620 & 4.44 & 15.97 \\
MUFusion & 0.574 & 0.229 & 0.189 & 1.196 & 0.604 & 36.75 & 0.020 & 0.429 & 0.589 & $-$0.440 & 0.083 & 0.343 & 0.387 & 0.679 & $-$1949 & 0.358 & 5.30 & 14.76 \\
NestFuse & 0.596 & 0.269 & 0.161 & 1.052 & 0.724 & 33.06 & 0.021 & 0.369 & 0.618 & $-$0.283 & 0.100 & 0.254 & 0.293 & 0.814 & $-$3497 & 0.655 & 4.34 & 14.88 \\
DIDFuse & 0.590 & 0.273 & 0.190 & 0.884 & 0.745 & 46.07 & 0.024 & 0.427 & 0.624 & $-$0.289 & 0.110 & 0.297 & 0.328 & 0.773 & $-$1916 & 0.345 & 4.67 & 14.74 \\
\bottomrule
\end{tabular}}
\end{table}

% ────────────────────────────────────────────────────────────────
%  TABLE — PET-MRI
% ────────────────────────────────────────────────────────────────
\begin{table}[H]
\centering
\caption{Quantitative comparison on \textbf{PET-MRI} (24 pairs).}
\label{tab:petmri}
\setlength{\tabcolsep}{2.5pt}
\renewcommand{\arraystretch}{1.1}
\scriptsize
\resizebox{\textwidth}{!}{%
\begin{tabular}{l | cc | cc | ccc | cccc | ccc | cc | cc}
\toprule
 &
\multicolumn{2}{c|}{\textbf{Similarity}} &
\multicolumn{2}{c|}{\textbf{Fidelity}} &
\multicolumn{3}{c|}{\textbf{Info Theory}} &
\multicolumn{4}{c|}{\textbf{Image Feature}} &
\multicolumn{3}{c|}{\textbf{Struct Sim}} &
\multicolumn{2}{c|}{\textbf{Perception}} &
\\
\cmidrule(lr){2-3}\cmidrule(lr){4-5}\cmidrule(lr){6-8}\cmidrule(lr){9-12}\cmidrule(lr){13-15}\cmidrule(lr){16-17}
\textbf{Method} &
\rot{MS-SSIM} & \rot{FMI} & \rot{VIF} & \rot{SCD} &
\rot{$Q_{MI}$} & \rot{$Q_{TE}$} & \rot{$Q_{NCIE}$} &
\rot{$Q^{AB/F}$} & \rot{$Q_M$} & \rot{$Q_{SF}$} & \rot{$Q_P$} &
\rot{$Q_S$} & \rot{$Q_C$} & \rot{$Q_Y$} &
\rot{$Q_{CV}$} & \rot{$Q_{CB}$} & \rot{EN} & \rot{PSNR} \\
\midrule
CDDFuse\textsubscript{IVF} & 0.674 & 0.305 & 0.328 & \best{1.815} & 0.745 & 93.05 & 0.021 & 0.646 & 0.699 & $-$0.071 & 0.157 & 0.371 & 0.392 & 0.947 & $-$764 & 0.678 & 4.15 & 14.08 \\
CDDFuse\textsubscript{MIF} & 0.691 & \best{0.350} & 0.353 & 1.683 & 0.804 & \best{133.56} & 0.025 & 0.697 & \best{0.752} & $-$0.021 & 0.223 & 0.395 & 0.412 & \best{0.955} & $-$324 & \best{0.743} & 4.14 & 15.97 \\
MambaDFuse & 0.689 & 0.310 & \best{0.381} & 1.529 & 0.659 & 21.36 & 0.038 & \best{0.747} & 0.704 & $-$0.047 & \best{0.327} & 0.533 & 0.563 & 0.604 & \best{$-$72} & 0.383 & \best{6.39} & 15.68 \\
SwinFusion & 0.685 & 0.214 & 0.273 & 1.244 & 0.535 & 8.00 & 0.019 & 0.667 & 0.528 & $-$0.156 & 0.176 & 0.512 & 0.524 & 0.561 & $-$435 & 0.363 & 5.99 & 14.72 \\
LRFNet & \best{0.700} & 0.281 & 0.347 & 1.456 & 0.727 & 18.36 & 0.033 & 0.703 & 0.615 & $-$0.152 & 0.301 & 0.522 & 0.545 & 0.936 & $-$112 & 0.607 & 5.30 & \best{16.87} \\
IFCNN & 0.694 & 0.240 & 0.275 & 1.221 & 0.601 & 9.17 & 0.022 & 0.667 & 0.568 & \best{$-$0.003} & 0.250 & \best{0.576} & \best{0.589} & 0.922 & $-$185 & 0.612 & 5.36 & 16.26 \\
EMMA & 0.688 & 0.228 & 0.296 & 1.577 & 0.593 & 11.08 & 0.027 & 0.667 & 0.537 & $-$0.123 & 0.178 & 0.525 & 0.547 & 0.908 & $-$107 & 0.498 & 6.22 & 15.77 \\
DDFM & 0.696 & 0.271 & 0.344 & 1.297 & \best{0.815} & 26.04 & \best{0.041} & 0.545 & 0.533 & $-$0.407 & 0.293 & 0.477 & 0.481 & 0.837 & $-$996 & 0.607 & 5.40 & 15.52 \\
MUFusion & 0.651 & 0.224 & 0.294 & 0.817 & 0.571 & 10.05 & 0.024 & 0.651 & 0.552 & $-$0.179 & 0.198 & 0.536 & 0.552 & 0.779 & $-$178 & 0.405 & 6.02 & 15.87 \\
NestFuse & 0.638 & 0.184 & 0.209 & 1.450 & 0.746 & 13.56 & 0.022 & 0.233 & 0.514 & $-$0.390 & 0.097 & 0.224 & 0.244 & 0.682 & $-$2608 & 0.567 & 4.06 & 14.26 \\
DIDFuse & 0.675 & 0.198 & 0.189 & 1.247 & 0.663 & 10.76 & 0.023 & 0.341 & 0.518 & $-$0.247 & 0.142 & 0.291 & 0.299 & 0.697 & $-$1359 & 0.359 & 4.86 & 14.61 \\
\bottomrule
\end{tabular}}
\end{table}

% ────────────────────────────────────────────────────────────────
%  TABLE — SPECT-MRI
% ────────────────────────────────────────────────────────────────
\begin{table}[H]
\centering
\caption{Quantitative comparison on \textbf{SPECT-MRI} (24 pairs).}
\label{tab:spectmri}
\setlength{\tabcolsep}{2.5pt}
\renewcommand{\arraystretch}{1.1}
\scriptsize
\resizebox{\textwidth}{!}{%
\begin{tabular}{l | cc | cc | ccc | cccc | ccc | cc | cc}
\toprule
 &
\multicolumn{2}{c|}{\textbf{Similarity}} &
\multicolumn{2}{c|}{\textbf{Fidelity}} &
\multicolumn{3}{c|}{\textbf{Info Theory}} &
\multicolumn{4}{c|}{\textbf{Image Feature}} &
\multicolumn{3}{c|}{\textbf{Struct Sim}} &
\multicolumn{2}{c|}{\textbf{Perception}} &
\\
\cmidrule(lr){2-3}\cmidrule(lr){4-5}\cmidrule(lr){6-8}\cmidrule(lr){9-12}\cmidrule(lr){13-15}\cmidrule(lr){16-17}
\textbf{Method} &
\rot{MS-SSIM} & \rot{FMI} & \rot{VIF} & \rot{SCD} &
\rot{$Q_{MI}$} & \rot{$Q_{TE}$} & \rot{$Q_{NCIE}$} &
\rot{$Q^{AB/F}$} & \rot{$Q_M$} & \rot{$Q_{SF}$} & \rot{$Q_P$} &
\rot{$Q_S$} & \rot{$Q_C$} & \rot{$Q_Y$} &
\rot{$Q_{CV}$} & \rot{$Q_{CB}$} & \rot{EN} & \rot{PSNR} \\
\midrule
CDDFuse\textsubscript{IVF} & \best{0.758} & 0.343 & 0.324 & \best{1.873} & 0.803 & $-$6.37 & 0.022 & 0.660 & 0.770 & $-$0.021 & 0.187 & 0.339 & 0.352 & 0.950 & $-$199 & 0.667 & 3.82 & 17.72 \\
CDDFuse\textsubscript{MIF} & 0.744 & \best{0.467} & \best{0.484} & 1.335 & \best{1.054} & $-$31.62 & 0.038 & \best{0.749} & \best{0.883} & \best{$-$0.010} & 0.326 & 0.380 & 0.423 & \best{0.990} & \best{$-$55} & \best{0.765} & 3.82 & \best{25.73} \\
MambaDFuse & 0.731 & 0.275 & 0.267 & 1.297 & 0.686 & 25.05 & 0.025 & 0.596 & 0.631 & $-$0.114 & 0.177 & 0.434 & 0.454 & 0.517 & $-$186 & 0.305 & 5.25 & 18.16 \\
SwinFusion & 0.722 & 0.253 & 0.248 & 0.832 & 0.625 & 21.85 & 0.022 & 0.591 & 0.606 & $-$0.108 & 0.154 & 0.431 & 0.454 & 0.518 & $-$200 & 0.329 & 5.44 & 18.93 \\
LRFNet & 0.724 & 0.334 & 0.327 & 1.375 & 0.809 & \best{35.93} & 0.032 & 0.655 & 0.710 & $-$0.144 & 0.292 & 0.434 & 0.468 & 0.954 & $-$61 & 0.650 & 4.71 & 22.27 \\
IFCNN & 0.727 & 0.300 & 0.288 & 1.004 & 0.699 & 24.90 & 0.026 & 0.631 & 0.677 & $-$0.017 & 0.264 & 0.561 & 0.583 & 0.945 & $-$111 & 0.654 & 4.90 & 21.22 \\
EMMA & 0.709 & 0.271 & 0.301 & 1.573 & 0.693 & 28.75 & 0.030 & 0.628 & 0.617 & $-$0.033 & 0.157 & 0.460 & 0.493 & 0.919 & $-$87 & 0.533 & \best{5.61} & 20.03 \\
DDFM & 0.716 & 0.274 & 0.273 & 1.794 & 0.788 & 13.38 & \best{0.041} & 0.542 & 0.579 & $-$0.237 & \best{0.338} & \best{0.621} & \best{0.624} & 0.867 & $-$287 & 0.536 & 5.60 & 17.90 \\
MUFusion & 0.681 & 0.277 & 0.277 & 1.096 & 0.691 & 24.10 & 0.028 & 0.586 & 0.638 & $-$0.183 & 0.217 & 0.482 & 0.515 & 0.787 & $-$132 & 0.387 & 5.27 & 21.35 \\
NestFuse & 0.676 & 0.212 & 0.165 & 1.304 & 0.655 & 17.80 & 0.016 & 0.345 & 0.583 & $-$0.082 & 0.084 & 0.249 & 0.259 & 0.791 & $-$855 & 0.611 & 3.68 & 17.07 \\
DIDFuse & 0.680 & 0.237 & 0.167 & 0.186 & 0.708 & 22.73 & 0.021 & 0.306 & 0.598 & $-$0.280 & 0.125 & 0.233 & 0.244 & 0.729 & $-$681 & 0.308 & 4.25 & 18.00 \\
\bottomrule
\end{tabular}}
\end{table}

\restoregeometry

% ════════════════════════════════════════════════════════════════
\begin{thebibliography}{99}
\bibitem{fusionmamba2024} X.~Xie et al., ``FusionMamba: Efficient image fusion with state space model,'' \textit{arXiv:2404.09498}, 2024.
\bibitem{liu2012} Z.~Liu et al., ``Objective assessment of multiresolution image fusion algorithms for context enhancement in night vision: A comparative study,'' \textit{IEEE TPAMI}, vol.~34, no.~1, pp.~94--109, 2012.
\bibitem{wang2003} Z.~Wang, E.~Simoncelli, and A.~Bovik, ``Multi-scale structural similarity for image quality assessment,'' \textit{Proc.\ IEEE Asilomar}, 2003.
\bibitem{haghighat2011} M.~Haghighat, M.~Razian, et al., ``A non-reference image fusion metric based on mutual information of image features,'' \textit{Comput.\ Electr.\ Eng.}, vol.~37, no.~5, 2011.
\bibitem{sheikh2006} H.~Sheikh and A.~Bovik, ``Image information and visual quality,'' \textit{IEEE TIP}, vol.~15, no.~2, pp.~430--444, 2006.
\bibitem{aslantas2015} V.~Aslantas and E.~Bendes, ``A new image quality metric for image fusion: The sum of the correlations of differences,'' \textit{AEU}, vol.~69, no.~12, 2015.
\bibitem{hossny2008} M.~Hossny, S.~Nahavandi, and D.~Creighton, ``Comments on `Information measure for performance of image fusion','' \textit{Electron.\ Lett.}, vol.~44, no.~18, 2008.
\bibitem{cvejic2006} N.~Cvejic, C.~Canagarajah, and D.~Bull, ``Image fusion metric based on mutual information and Tsallis entropy,'' \textit{Electron.\ Lett.}, vol.~42, no.~11, 2006.
\bibitem{nava2007} R.~Nava, G.~Crist\'{o}bal, and B.~Escalante-Ram\'{i}rez, ``Mutual information improves image fusion quality assessments,'' \textit{SPIE News Room}, 2007.
\bibitem{xydeas2000} C.~Xydeas and V.~Petrovic, ``Objective image fusion performance measure,'' \textit{Electron.\ Lett.}, vol.~36, no.~4, pp.~308--309, 2000.
\bibitem{wang2008} P.~Wang and B.~Liu, ``A novel image fusion metric based on multi-scale analysis,'' \textit{Proc.\ IEEE ICASSP}, pp.~965--968, 2008.
\bibitem{zheng2007} Y.~Zheng et al., ``A new metric based on extended spatial frequency and its application to DWT based fusion algorithms,'' \textit{Inf.\ Fusion}, vol.~8, no.~2, 2007.
\bibitem{zhao2007} J.~Zhao, R.~Lagani\`{e}re, and Z.~Liu, ``Performance assessment of combinative pixel-level image fusion based on an absolute feature measurement,'' \textit{IJICIC}, vol.~3, no.~6(A), 2007.
\bibitem{piella2003} G.~Piella and H.~Heijmans, ``A new quality metric for image fusion,'' \textit{Proc.\ ICIP}, 2003.
\bibitem{cvejic2005} N.~Cvejic, A.~Loza, D.~Bull, and N.~Canagarajah, ``A similarity metric for assessment of image fusion algorithms,'' \textit{Int'l J.\ Signal Process.}, vol.~2, no.~3, 2005.
\bibitem{yang2008} C.~Yang, J.~Zhang, X.~Wang, and X.~Liu, ``A novel similarity based quality metric for image fusion,'' \textit{Inf.\ Fusion}, vol.~9, pp.~156--160, 2008.
\bibitem{chen2007} H.~Chen and P.~Varshney, ``A human perception inspired quality metric for image fusion based on regional information,'' \textit{Inf.\ Fusion}, vol.~8, pp.~193--207, 2007.
\bibitem{chen2009} Y.~Chen and R.~Blum, ``A new automated quality assessment algorithm for image fusion,'' \textit{Image Vis.\ Comput.}, vol.~27, pp.~1421--1432, 2009.
\bibitem{cddfuse} J.~Zhao et al., ``CDDFuse: Correlation-driven dual-branch feature decomposition for multi-modality image fusion,'' \textit{CVPR}, 2023.
\bibitem{mambadfuse} Z.~Li et al., ``MambaDFuse: A Mamba-based dual-phase model for multi-modality image fusion,'' \textit{arXiv:2404.08406}, 2024.
\bibitem{swinfusion} J.~Ma et al., ``SwinFusion: Cross-domain long-range learning for general image fusion via Swin Transformer,'' \textit{IEEE/CAA JAS}, vol.~9, no.~7, 2022.
\bibitem{lrfnet} D.~He et al., ``LRFNet: A real-time medical image fusion method guided by detail information,'' 2023.
\bibitem{ifcnn} Y.~Zhang et al., ``IFCNN: A general image fusion framework based on convolutional neural network,'' \textit{Inf.\ Fusion}, vol.~54, 2020.
\bibitem{emma} L.~Cao et al., ``EMMA: Efficient multi-modal adapter for medical image fusion,'' 2024.
\bibitem{ddfm} Z.~Zhao et al., ``DDFM: Denoising diffusion model for multi-modality image fusion,'' \textit{ICCV}, 2023.
\bibitem{mufusion} Y.~Cheng et al., ``MUFusion: A general unsupervised image fusion network based on memory unit,'' \textit{Inf.\ Fusion}, vol.~92, 2023.
\bibitem{nestfuse} H.~Li et al., ``NestFuse: An infrared and visible image fusion architecture based on nest connection and spatial/channel attention models,'' \textit{IEEE TIM}, vol.~69, no.~12, 2020.
\bibitem{didfuse} J.~Zhao et al., ``DIDFuse: Deep image decomposition for infrared and visible image fusion,'' \textit{IJCAI}, 2020.
\end{thebibliography}

\end{document}
